\begin{proofbox}
\textbf{证明前提}

设$p>0$，$P(x_0,y_0)$在抛物线$y^2=2px$上，因此
\[
    y_0^2=2px_0.
\]
下面按\textbf{最低门槛 $\rightarrow$ 证明最快 $\rightarrow$ 通用证明}排列。
三种方法最终都要得到$\boxed{y_0y=p(x+x_0)}$，而且都包含顶点$P=O$的情形。

\medskip
\begin{minipage}{\linewidth}
\textbf{方法一：最低门槛——两点相减与平方差（不使用导数）}
\begin{center}\begin{tikzpicture}[>=stealth,font=\small]
\node[anchor=west,font=\small\bfseries] at (-.6,4.95) {平方差法：割线方向由坐标增量确定};
\draw[->,gray] (-.6,0)--(4.2,0) node[right] {$x$};
\draw[->,gray] (0,-1.2)--(0,4.1) node[above] {$y$};
\draw[blue,very thick,domain=-1.1:3.9,samples=100] plot ({\x*\x/4},\x);
\node[blue,fill=white] at (3.3,3.05) {$y^2=4x$};
\node[below left] at (0,0) {$O$};
\draw[orange!80!black,very thick] (-.5,.5)--(3,4);
\draw[teal,dashed,thick] (-.5,1)--(3.9,3.9333);
\coordinate (P) at (1,2); \coordinate (Q) at (2.25,3);
\draw[gray,densely dotted,thick] (P)--(2.25,2)--(Q);
\node[below] at (1.65,2) {$x_1-x_0$};
\node[right] at (2.25,2.5) {$y_1-y_0$};
\fill (P) circle (2pt);\node[anchor=east] at (-.5,1.7) {$P(x_0,y_0)$};\draw[gray] (-.4,1.7)--(.85,1.95);
\fill (Q) circle (2pt) node[above left=4pt] {$Q(x_1,y_1)$};
\node[anchor=west] at (4.5,3.5) {虚线：割线 $PQ$};
\node[anchor=west,orange!80!black] at (4.5,2.6) {实线：$P$ 处切线};
\node[anchor=west] at (4.5,1.4) {$\dfrac{x_1-x_0}{y_1-y_0}=\dfrac{y_1+y_0}{2p}$};
\node[anchor=west] at (4.5,.2) {$Q\to P$ 时，比值 $\to y_0/p$};
\end{tikzpicture}
\end{center}
\end{minipage}\par\medskip

\textbf{第一步：}在抛物线上另取$Q(x_1,y_1)\ne P$。两个点的方程相减：
\[
 y_1^2-y_0^2=2p(x_1-x_0)
 \quad\Longrightarrow\quad
 (y_1-y_0)(y_1+y_0)=2p(x_1-x_0).
\]
\textbf{为什么相减？}相减后出现$x_1-x_0$、$y_1-y_0$，正好是从$P$到$Q$的坐标变化量，
可以确定割线$PQ$的方向。

\textbf{第二步：}由于同一抛物线上$y$相同的两点必重合，故$Q\ne P$时$y_1\ne y_0$，于是
\[
    \frac{x_1-x_0}{y_1-y_0}=\frac{y_1+y_0}{2p}.
\]
当$Q$沿抛物线越来越靠近$P$时，$y_1$越来越接近$y_0$，
割线的方向越来越接近切线的方向，得到切线方程
\[
    x-x_0=\frac{y_0}{p}(y-y_0).
\]
这里写$x$相对$y$的变化，\textbf{不用写$\frac{p}{y_0}$}，所以顶点处也适用。

\textbf{第三步：}两边乘以$p$并使用$y_0^2=2px_0$：
\[
    px-px_0=y_0y-y_0^2=y_0y-2px_0
    \quad\Longrightarrow\quad
    \boxed{y_0y=p(x+x_0)}.
\]
\textbf{不靠图形的检查：}把该直线写为$x=\dfrac{y_0}{p}y-x_0$，代回抛物线得到
\[
    y^2=2y_0y-2px_0=2y_0y-y_0^2
    \quad\Longrightarrow\quad (y-y_0)^2=0.
\]
这说明\textbf{两个交点对应的解重合在$P$}，与割线变成切线的结果一致。

\medskip
\begin{minipage}{\linewidth}
\textbf{方法二：证明最快——直接求导（两步）}
\begin{center}\begin{tikzpicture}[>=stealth,font=\small]
\node[anchor=west,font=\small\bfseries] at (-.6,4.95) {导数法：读出切线上 $x$ 随 $y$ 的变化};
\draw[->,gray] (-.6,0)--(4.2,0) node[right] {$x$};
\draw[->,gray] (0,-1.2)--(0,4.1) node[above] {$y$};
\draw[blue,very thick,domain=-1.1:3.9,samples=100] plot ({\x*\x/4},\x);
\node[blue,fill=white] at (3.3,3.05) {$y^2=4x$};
\node[below left] at (0,0) {$O$};
\draw[orange!80!black,very thick] (-.5,.5)--(3,4);
\draw[teal,thick,dashed] (1,2)--(2,2)--(2,3);
\fill (1,2) circle (2pt) node[above left=4pt] {$P$};
\fill (2,3) circle (2pt) node[above left=4pt] {$R$};
\node[below] at (1.5,2) {$\Delta x$};
\node[right] at (2,2.5) {$\Delta y$};
\node[anchor=west] at (4.5,3.6) {$R$ 在切线上，$\Delta y\ne0$};
\node[anchor=west] at (4.5,2.5) {$\dfrac{\Delta x}{\Delta y}=\left.\dfrac{\mathrm dx}{\mathrm dy}\right|_P=\dfrac{y_0}{p}$};
\node[anchor=west] at (4.5,1.2) {顶点处：$\Delta x=0$};
\node[anchor=west] at (4.5,.3) {故切线竖直，$x=0$};
\end{tikzpicture}
\end{center}
\end{minipage}\par\medskip

\textbf{第一步：}将抛物线写为$x=\dfrac{y^2}{2p}$，\textbf{把$y$看作自变量}，求导得
\[
    f'(y)=\frac{y}{p}.
\]
所以在$P$处，切线可写成
\[
    x-x_0=\frac{y_0}{p}(y-y_0).
\]
\textbf{第二步：}乘以$p$，再用$y_0^2=2px_0$：
\[
    p(x-x_0)=y_0(y-y_0)
    \quad\Longrightarrow\quad
    \boxed{y_0y=p(x+x_0)}.
\]
\textbf{顶点检查：}$y_0=x_0=0$时，上式直接成为$x=0$。
因为这里把$y$当作自变量，对函数$x=f(y)$求导，
无需除以$y_0$，所以不会漏掉顶点切线。

\medskip
\begin{minipage}{\linewidth}
\textbf{方法三：通用证明——设直线并用判别式}
\begin{center}\begin{tikzpicture}[>=stealth,font=\small]
\foreach \shift/\b/\title in {0/0/割线：两个不同交点,6.2/-1/切线：两个根重合}{
\begin{scope}[xshift=\shift cm]
\node[anchor=west,font=\small\bfseries] at (-.5,4.8) {\title};
\draw[->,gray] (-.5,0)--(4.5,0) node[right] {$x$};
\draw[->,gray] (0,-.8)--(0,4.6);\node[left] at (0,3.3) {$y$};
\draw[blue,very thick,domain=-.7:4.2,samples=100] plot ({\x*\x/4},\x);
\draw[orange!80!black,very thick] ({-.5+\b},-.5)--({4.3+\b},4.3);
\node at (2,-1.4) {\ifnum\b=0 $x=y$\else $x=y-1$\fi};
\end{scope}}
\fill (0,0) circle (2pt) node[below right=6pt] {$O$};
\fill (4,4) circle (2pt) node[below right=6pt] {$(4,4)$};
\fill (7.2,2) circle (2pt) node[below right=7pt] {$P(1,2)$};
\node at (2,-2.2) {$y(y-4)=0$};
\node at (8.2,-2.2) {$(y-2)^2=0$，$\Delta=0$};
\end{tikzpicture}
\end{center}
\end{minipage}\par\medskip

\textbf{第一步：}由于本抛物线的切线没有水平方向，统一设直线为
\[
    x=my+b.
\]
这样的写法能包含\textbf{竖直切线$x=0$}。代入$y^2=2px$：
\[
    y^2-2pmy-2pb=0.
\]
\textbf{第二步：}当直线与抛物线相切时，以上关于$y$的一元二次方程有两个相等的实数根，
所以判别式
\[
    \Delta=(-2pm)^2-4(1)(-2pb)=4p^2m^2+8pb=0,
\]
得到$b=-\dfrac{pm^2}{2}$。方程配方为
\[
    (y-pm)^2=0,
\]
故切点纵坐标$y_0=pm$，从而$m=\dfrac{y_0}{p}$。
由$y_0^2=2px_0$知$b=-\dfrac{y_0^2}{2p}=-x_0$。

\textbf{第三步：}代回直线方程
\[
    x=\frac{y_0}{p}y-x_0
    \quad\Longrightarrow\quad
    \boxed{y_0y=p(x+x_0)}.
\]
\textbf{适用提醒：}若一开始设成$y=kx+b$，就会漏掉顶点处的竖直切线；
若遇到水平直线$y=c$，即使只有一个公共点，也\textbf{不能}把它当成切线。

\medskip
\textbf{推论证明：轴上截距等于切点横坐标的相反数}

\textbf{方法甲（最快，直接代入）：}设切线交$x$轴于$T(t,0)$，
把$y=0$代入$y_0y=p(x+x_0)$，得
\[
  0=p(t+x_0)\quad\Longrightarrow\quad\boxed{t=-x_0}.
\]
\textbf{方法乙（最低门槛，观察直角三角形）：}当$P\ne O$时，$x_0>0$，
切线斜率$k=\dfrac{p}{y_0}$。又因$P$在抛物线上，
$y_0^2=2px_0$，故
\[
  k=\frac{p}{y_0}=\frac{y_0}{2x_0}.
\]
设切线交$x$轴于$T(t,0)$。由切线过$P(x_0,y_0)$，
其斜率也等于$\dfrac{y_0}{x_0-t}$，于是
\[
  \frac{y_0}{x_0-t}=\frac{y_0}{2x_0}
  \quad\Longrightarrow\quad x_0-t=2x_0
  \quad\Longrightarrow\quad\boxed{t=-x_0}.
\]
过$P$作$PH\perp x$轴，$H(x_0,0)$。于是
\[
  T(-x_0,0),\qquad O(0,0),\qquad H(x_0,0),
  \qquad TO=OH=x_0.
\]
由此得到\textbf{$O$是$TH$的中点}。
当$P=O$时切线为$x=0$，$T=H=O$，截距结论仍成立。

\textbf{与通用证明的联系：}方法三已推出直线$x=my+b$的$b=-x_0$；
令$y=0$就有$x=b=-x_0$，因此三种证明都能顺手得到新推论。

\medskip
\textbf{三法对照}

最低门槛法看的是\textbf{平方差给出割线方向}；最快法看的是\textbf{导数给出切线方向}；
通用法看的是\textbf{二次方程出现相等的根}。三法结论相同，但入手信号不同。
\end{proofbox}
